\section{引言}

碳化硅（SiC）作为第三代宽禁带半导体材料，在新能源汽车、智能电网、5G 通信等领域的功率器件中具有不可替代的地位。外延层是 SiC 功率器件的核心功能区，其厚度直接影响器件的击穿电压和导通电阻。以 1200V 级 SiC MOSFET 为例，漂移层厚度约为 10\,$\mu$m，厚度偏差 0.5\,$\mu$m 即可导致击穿电压偏移数十伏特\cite{ref-s15}。因此，外延层厚度的精确、无损测量是 SiC 器件制造中不可或缺的质量控制环节。

红外反射法（Infrared Reflectance Spectroscopy）是半导体工业中测量外延层厚度的标准方法之一，已被纳入 SEMI MF95\cite{ref-s01} 和中国国家标准 GB/T 42905—2023\cite{ref-s02}。其基本原理是：红外光在外延层上下界面反射后发生干涉，反射率谱呈现周期性条纹，条纹频率与光在外延层中的往返光程成正比，从而可由条纹周期反演厚度。该方法无需破坏样品、无需标定参照，在工业产线中广泛使用。

然而，该方法的实际应用面临若干挑战：（1）折射率的色散与厚度存在简并关系，单角度数据难以同时辨识；（2）衬底界面的反射相位并非理想常数，重掺衬底的 Drude 自由载流子效应使相位随频率变化；（3）当外延层与衬底的折射率差较大时，高阶反射束（多光束干涉）不可忽略，双光束近似失效。这些因素共同构成了厚度反演的系统误差来源。

本文针对 2025 年全国大学生数学建模竞赛 B 题，以红外反射光谱数据为输入，依次解决三个问题：（1）建立双光束干涉数学模型，导出厚度与反射率谱的定量关系；（2）设计多算法融合的厚度反演框架，给出 SiC 外延层厚度的最佳估计及其不确定度；（3）将模型升级为 Airy 多光束形式，判断 Si 和 SiC 样品中多光束干涉的显著性，并量化其对厚度估计的影响。

本文的模型体系以"判据先行、实验拍板"为方法论核心：每个建模决策（折射率取值、相位模型、偏振处理等）均以预注册的定量判据经数据实验定案，落选方案的量化结果全部归档，作为系统误差区间和模型升级的依据。

\subsection{问题重述}

\textbf{问题1：} 建立双光束干涉数学模型。根据附件1（10\degree 入射角）和附件2（15\degree 入射角）的红外反射率谱数据，导出反射率与波数、外延层厚度之间的定量关系，确定条纹级数与厚度的线性关系，并讨论模型的可解性与唯一性。

\textbf{问题2：} 设计多算法融合的厚度反演算法。利用双光束模型从附件1/2 数据中反演 SiC 外延层厚度，给出最终估计值及其统计不确定度与系统误差区间，并对结果的可靠性进行多维度验证。

\textbf{问题3：} 多光束干涉效应的判定与修正。考虑外延层内多次往返反射的相干叠加（Airy 多光束模型），推导多光束干涉可观测的必要条件；分别判断 Si 样品（附件3/4）和 SiC 样品（附件1/2）中多光束干涉的显著性；对 Si 样品给出厚度反演结果，对 SiC 样品给出双光束模型的修正量。